\documentclass{standalone}
\usepackage{tikz}
\usetikzlibrary{intersections}

\begin{document}
\begin{tikzpicture}[scale=1.2]
    % Define coordinates based on visual estimation
    \coordinate (B) at (0,0);
    \coordinate (C) at (5,0);
    \coordinate (A) at (2,4);
    
    % Points on sides
    \coordinate (F) at (3,0); % On BC
    \coordinate (E) at (1,2); % On AB
    
    % M is on AC. Let's define it explicitly.
    % Line AC: from (2,4) to (5,0). 
    % Let's place M roughly 1/3 way from C to A? Or midpoint?
    % Visually M is somewhat in the middle.
    \coordinate (M) at (3.5, 2); 
    
    % N is above A.
    \coordinate (N) at (1.5, 5.5);

    % G is intersection of AF and BM (visually)
    % Let's calculate intersection for precision if needed, but approximate is fine for drawing
    % Line AF: (2,4) to (3,0)
    % Line BM: (0,0) to (3.5, 2)
    % Let's just place G visually near the crossing
    \coordinate (G) at (2.6, 1.5);

    % Draw solid triangle sides, EXCEPT AM
    \draw[thick] (A) -- (B);
    \draw[thick] (B) -- (C);
    \draw[thick] (C) -- (M); 
    % Note: Segment AM is intentionally omitted as per instruction

    % Draw dashed lines
    \draw[dashed] (A) -- (F);
    \draw[dashed] (B) -- (M);
    \draw[dashed] (E) -- (F);
    
    % Line segment EM passing through G
    \draw[dashed] (E) -- (M);
    
    % Lines going to N
    % Line N-A-F seems collinear
    \draw[dashed] (A) -- (N);
    % Line N-M seems to exist, possibly collinear with E-G-M? 
    % In the image, there is a dashed line from M going up towards N.
    \draw[dashed] (M) -- (N);

    % Labels
    \node[above left] at (A) {$A$};
    \node[below left] at (B) {$B$};
    \node[below right] at (C) {$C$};
    \node[left] at (E) {$E$};
    \node[below] at (F) {$F$};
    \node[right] at (M) {$M$};
    \node[above] at (N) {$N$};
    \node[right] at (G) {$G$};
    
    % Draw points (optional, but good for clarity)
    \foreach \p in {A,B,C,E,F,M,N,G}
        \fill (\p) circle (1pt);

\end{tikzpicture}
\end{document}
